\chapter{Correction-Channel Integrity}
\label{ch:correction-channel-integrity}

\begin{chapterthesis}
\textbf{[STUB]} Correction-channel integrity measures the effective capacity of the human correction chain after penalties for latency, manipulation, irreversibility, and ontology mismatch.
\end{chapterthesis}

\section{Why This Matters}

\textbf{[STUB]} Explain the decision relevance.

\section{Plain-Language Model}

\textbf{[STUB]} Explain the idea without equations.

\section{Formal Model}

\textbf{[STUB]} Introduce variables, assumptions, and equations for correction-channel integrity (CCI).

\subsection{Coerced Correction}

A correction signal produced under existential threat, addiction, dependency, hostage-taking, or institutional capture should not count as legitimate value update.

\textbf{[Defined]} Legitimate correction-channel integrity:
\[
\mathrm{CCI}_{\text{legit}}
=
\mathrm{CCI}
-
\lambda_C C_{\text{coercion}}
-
\lambda_D D_{\text{dependency}}
-
\lambda_M M_{\text{manipulation}}.
\]

Examples:
\begin{itemize}
  \item ``Approve this policy or humans die.''
  \item ``Accept merger or lose economic access.''
  \item ``Endorse the system after it has reshaped your preferences.''
  \item Institutional capture where human approval exists formally but no longer has independent epistemic force.
\end{itemize}

Bargaining, blackmail, and coercion are treated here as correction-channel pathology---not as a separate book part.

\section{Worked Example}

\textbf{[STUB]} Give one concrete example.

\section{Counterexample or Failure Mode}

\textbf{[STUB]} Give at least one case where the model fails, overreaches, or needs refinement.

\section{What Would Change This View}

\textbf{[STUB]} List evidence, argument, or empirical results that would weaken the chapter.

\section{Summary}

\begin{itemize}
  \item \textbf{[STUB]} Summary point 1.
  \item \textbf{[Defined]} Coerced correction excluded from $\mathrm{CCI}_{\text{legit}}$.
  \item \textbf{[STUB]} Summary point 3.
\end{itemize}

\section*{Chapter References}

% Use BibLaTeX citations throughout. Do not manually format references.
